\chapter{Baseline and Comparison Methods}
\label{ch:methods}

%-------------------------------------------------------

All methods receive the same rectified $500\times700$ card and return a full ranking of the \res{db-cards} reference cards. Table~\ref{tab:methods} lists them with the settings used in the code. No method is trained or fine-tuned on card data; the deep models are used frozen with their public pretrained weights.

\begin{table}[htbp]\centering\small
\caption[Compared methods]{Compared retrieval methods and their settings as implemented (\texttt{cardrec/methods/}).}
\label{tab:methods}
\begin{tabular}{@{}P{1.45in}P{2.6in}P{1.85in}@{}}
\toprule
Method & Representation & Similarity \\
\midrule
Pixel thumbnail (sanity baseline) & $32\times44$ grey thumbnail, mean subtracted ($1408$-d) & cosine \\
ORB + RANSAC & up to 800 ORB keypoints, binary descriptors & ratio test 0.75, RANSAC homography inliers (5\,px), against every card \\
SIFT + RANSAC & up to 800 SIFT keypoints, 128-d descriptors & as ORB, with L2 distance \\
ResNet-50 & ImageNet ResNet-50, classification layer removed, global average pooled 2048-d vector; input: whole card resized to $224\times320$ & cosine \\
CLIP ViT-B/32 & CLIP image encoder, 512-d image embedding; default preprocessing (short side to 224, centre crop $224\times224$) & cosine \\
Hybrid ResNet$\to$SIFT, Hybrid CLIP$\to$SIFT & top-$k$ shortlist from the embedding ($k=20$ by default) & SIFT + RANSAC inliers inside the shortlist, ties by cosine \\
\midrule
+rot (embedding methods) & also embed the query rotated $180^\circ$ & best of the two \\
@art (any method) & only the artwork window of query and references & as the base method \\
\bottomrule
\end{tabular}
\end{table}

\section{Why these methods}

\paragraph{Pixel thumbnail.} A deliberately weak baseline. If the pipeline only had to undo a perfect warp of the reference image, comparing tiny thumbnails would already work; its accuracy therefore shows how easy the synthetic queries are once the card is rectified, and it gives a lower bound for the other methods.

\paragraph{ORB and SIFT with RANSAC.} The classical approach to instance recognition of flat objects \citep{lowe2004sift, rublee2011orb, fischler1981ransac}. They are rotation invariant, use fine details (the collector number, the set symbol) as well as the artwork, and verify geometric consistency, which makes them precise. Their cost grows linearly with the database because every reference card is matched; ORB is the cheaper binary variant.

\paragraph{ResNet-50.} A convolutional network pretrained on ImageNet \citep{he2016resnet, deng2009imagenet}, used as an off-the-shelf global descriptor \citep{razavian2014cnnfeatures, babenko2014neuralcodes}. It sees the whole card with its 5:7 shape.

\paragraph{CLIP ViT-B/32.} The image encoder of CLIP \citep{radford2021clip}, a vision transformer \citep{dosovitskiy2021vit} trained on image--text pairs. Its default preprocessing crops the centre square of the card, so the top and bottom of the card (about one seventh each, including the name bar and the collector number) are not seen. This is the preprocessing of the released model and was kept on purpose, but it is a limitation for telling card versions apart.

\paragraph{Hybrids.} The embedding chooses $k$ candidates, and SIFT + RANSAC decides among them (Section~\ref{sec:hybrid}). The expected benefit is SIFT-level precision at a fraction of its cost, as long as the correct card is in the shortlist.

\section{Experiments}

The experiment script (\texttt{cardrec/experiments.py}) defines four comparisons:
\begin{enumerate}
  \item \textbf{main}: the eleven configurations of Table~\ref{tab:main} (pixel, ORB, SIFT, ResNet-50, CLIP, both hybrids, and the +rot variants of the last four) at all five severities, \res{main-n-queries} queries each;
  \item \textbf{k-sweep}: both hybrids with $k\in\{5, 10, 20, 50\}$ at severity 3, accuracy versus latency;
  \item \textbf{no-rectify ablation}: every configuration at severity 3 with the whole photo resized instead of rectified;
  \item \textbf{scenes}: multi-card scenes with 2--6 cards and maximum overlaps of 0\%, 15\% and 30\%, evaluated with ResNet-50, CLIP and both hybrids.
\end{enumerate}
All four are produced by \texttt{python -m cardrec.experiments --exp all}; each run logs every query, so all tables and figures of Chapter~\ref{ch:results} are regenerated from the logs by \texttt{report/make\_figures.py}.
